% 一维教学例：位置0至7，A={1,2,3,6}，结构元B={-1,0,1}，区间外按背景处理。宽约 84mm。
\begin{tikzpicture}[x=5.5mm,y=5.5mm,font=\figurefont\small]
  \foreach \row/\name/\cells in {
    0/{$\bm{A}$}/{1,2,3,6},
    1/{$\bm{A}\ominus\bm{B}$}/{2},
    2/{$\bm{A}\oplus\bm{B}$}/{0,1,2,3,4,5,6,7},
    3/{$\bm{A}\circ\bm{B}$}/{1,2,3},
    4/{$\bm{A}\bullet\bm{B}$}/{1,2,3,4,5,6}}{
    \node[anchor=east] at (-.4,-\row*1.35+.5) {\name};
    \foreach \c in \cells{\fill[FigureBlue] (\c,-\row*1.35) rectangle ++(1,1);}
    \foreach \c in {0,...,7}{\draw[FigureStroke,line width=.4pt] (\c,-\row*1.35) rectangle ++(1,1);}
  }
  \foreach \c in {0,...,7}{\node[font=\figurefont\footnotesize,text=FigureMuted] at (\c+.5,1.5) {\c};}
  \begin{scope}[xshift=55mm]
    \node[anchor=east] at (-.4,.5) {$\bm{B}$};
    \foreach \c in {0,1,2}{\fill[FigureYellowFill] (\c,0) rectangle ++(1,1);
      \draw[FigureStroke,line width=.4pt] (\c,0) rectangle ++(1,1);}
    \fill[FigureStroke] (1.5,.5) circle (1.6pt);
    \node[font=\figurefont\footnotesize,text=FigureMuted,anchor=north] at (1.5,-.1) {原点};
  \end{scope}
\end{tikzpicture}
